\documentclass{zkdl-presentation-template}

\usepackage{changepage}

% Circuit notation of the book
\providecommand{\Circ}{\textit{\texttt{C}}}

% Plots
\pgfplotsset{
    zkdlgrid/.style={
        grid = both,
        grid style = {line width=.1pt, draw=oc-gray-2},
        major grid style = {line width=.2pt, draw=oc-gray-5}
    }
}

% Protocol diagrams
\tikzset{
    arrow/.style={-latex, ultra thick}
}

% Highlight of what changed in the protocol (as in the book)
\newcommand{\chg}[1]{\textcolor{oc-green-8}{#1}}

\title[Groth16]{\textbf{Pairing-Based SNARKs:\\Pinocchio and Groth16}}
\author{Blockstream Research}
\date{\today}
\homepage{research.blockstream.com}
\github{BlockstreamResearch/zk-book}

\begin{document}
    \maketitleframe

    \begin{frame}{Plan}
        \tableofcontents
    \end{frame}

    \section{Recap}

    \begin{frame}{Recap: R1CS}
        Each \textbf{constraint} in the Rank-1 Constraint System must be in the form:
        \begin{equation*}
            \langle \boldsymbol{\ell}, \mathbf{z}\rangle \times \langle \boldsymbol{r}, \mathbf{z}\rangle = \langle \boldsymbol{o}, \mathbf{z}\rangle
        \end{equation*}\pause

        Where $\langle \mathbf{u}, \mathbf{v}\rangle$ is a dot product.
        \vspace{-10pt}
        \begin{equation*}
            \langle \mathbf{u}, \mathbf{v} \rangle := \mathbf{u}^{\top}\mathbf{v} = \sum_{i=1}^{n} u_i v_i
            \vspace{-5pt}
        \end{equation*}\pause

        Thus
        \vspace{-5pt}
        \begin{equation*}
            \left(\sum_{i=1}^{n} \ell_i z_i\right) \times \left(\sum_{j=1}^{n} r_j z_j\right) = \sum_{k=1}^{n} o_k z_k
        \end{equation*}
        That is, actually, a quadratic equation with multiple variables.
    \end{frame}

    \begin{frame}[fragile]{Recap: R1CS}
        Consider the simplest program:
        \vspace{10pt}
        \begin{lstlisting}[language=Python,numbers=none]
    def example(a: F, b: F, c: F) -> F:
        if a:
            return b * c
        else:
            return b + c
        \end{lstlisting}
    \end{frame}

    \begin{frame}{Recap: R1CS}
        \vspace{-10pt}
        \begin{equation*}
            r = x_1 \times (x_2 \times x_3) + (1 - x_1) \times (x_2 + x_3)
        \end{equation*}\pause
        Thus, the next constraints can be built:
        \vspace{-5pt}
        \begin{align*}
            x_1 \times x_1 &= x_1 \quad \text{(binary check)} \tag{1} \\
            x_2 \times x_3 &= \mathsf{mult} \tag{2} \\
            x_1 \times \mathsf{mult} &= \mathsf{selectMult} \tag{3} \\
            (1 - x_1) \times (x_2 + x_3) &= r - \mathsf{selectMult} \tag{4}
        \end{align*}\pause
        The witness vector: $\mathbf{z} = (1, r, x_1, x_2, x_3, \mathsf{mult}, \mathsf{selectMult})$.

        \vspace{2pt}\pause
        The coefficients vectors:
        \vspace{-25pt}
        {\center\small\begin{align*}
            \boldsymbol{\ell}_1 &= (0, 0, 1, 0, 0, 0, 0), & \boldsymbol{r}_1 &= (0, 0, 1, 0, 0, 0, 0), & \boldsymbol{o}_1 &= (0, 0, 1, 0, 0, 0, 0) \\
            \boldsymbol{\ell}_2 &= (0, 0, 0, 1, 0, 0, 0), & \boldsymbol{r}_2 &= (0, 0, 0, 0, 1, 0, 0), & \boldsymbol{o}_2 &= (0, 0, 0, 0, 0, 1, 0) \\
            \boldsymbol{\ell}_3 &= (0, 0, 1, 0, 0, 0, 0), & \boldsymbol{r}_3 &= (0, 0, 0, 0, 0, 1, 0), & \boldsymbol{o}_3 &= (0, 0, 0, 0, 0, 0, 1) \\
            \boldsymbol{\ell}_4 &= (1, 0, -1, 0, 0, 0, 0), & \boldsymbol{r}_4 &= (0, 0, 0, 1, 1, 0, 0), & \boldsymbol{o}_4 &= (0, 1, 0, 0, 0, 0, -1)
        \end{align*}}
    \end{frame}

    \begin{frame}{Recap: QAP}
        R1CS provides us with the following constraint vectors:
        \vspace{-8pt}
        \begin{equation*}
            \boldsymbol{\ell}_1, \boldsymbol{\ell}_2, \dots, \boldsymbol{\ell}_m, \quad
            \boldsymbol{r}_1, \boldsymbol{r}_2, \dots, \boldsymbol{r}_m, \quad
            \boldsymbol{o}_1, \boldsymbol{o}_2, \dots, \boldsymbol{o}_m.
            \vspace{-5pt}
        \end{equation*}\pause
        Of course, they form corresponding matrices:
        \vspace{-5pt}
        \begin{equation*}
            L = \begin{bmatrix}
                L_{1,1} & L_{1,2} & \dots & L_{1,n} \\
                L_{2,1} & L_{2,2} & \dots & L_{2,n} \\
                \vdots & \vdots & \ddots & \vdots \\
                L_{m,1} & L_{m,2} & \dots & L_{m,n}
            \end{bmatrix}, \;
            \text{same goes for $R$ and $O$}
        \end{equation*}\pause
        An example of a single ``if'' statement:
        \begin{center}
            \begin{minipage}{0.4\textwidth}
            \vspace{-15pt}
            {\small \begin{align*}
                \boldsymbol{\ell}_1 &= (0, 0, 1, 0, 0, 0, 0) \\
                \boldsymbol{\ell}_2 &= (0, 0, 0, 1, 0, 0, 0) \\
                \boldsymbol{\ell}_3 &= (0, 0, 1, 0, 0, 0, 0) \\
                \boldsymbol{\ell}_4 &= (1, 0, -1, 0, 0, 0, 0)
            \end{align*}}
            \end{minipage}
            \begin{minipage}{0.5\textwidth}
                \vspace{-10pt}
                \begin{equation*}
                    \begin{bmatrix}
                        0 & 0 & \textcolor{oc-red-7}{1} & 0 & 0 & 0 & 0 \\
                        0 & 0 & \textcolor{oc-red-7}{0} & 1 & 0 & 0 & 0 \\
                        0 & 0 & \textcolor{oc-red-7}{1} & 0 & 0 & 0 & 0 \\
                        \textcolor{oc-blue-7}{1} & \textcolor{oc-blue-7}{0} & \boxed{-1} & \textcolor{oc-blue-7}{0} & \textcolor{oc-blue-7}{0} & \textcolor{oc-blue-7}{0} & \textcolor{oc-blue-7}{0}
                    \end{bmatrix}
                \end{equation*}
            \end{minipage}
        \end{center}
    \end{frame}

    \begin{frame}{Recap: QAP}
        \begin{center}
            \begin{tikzpicture}
                \begin{axis}[
                    axis lines = middle,
                    xlabel = {$T$},
                    ylabel = {$\ell_{x_1}(T)$},
                    ymin = -2.99, ymax = 2.99,
                    xmin = -0.5, xmax = 4.99,
                    domain = 0:5,
                    samples = 100,
                    ytick = {-3,...,3},
                    xtick = {0,1,...,5},
                    zkdlgrid
                ]
                \addplot[
                    color=oc-blue-7,
                    thick
                ]
                {-5/6*x^3 + 6*x^2 - 79/6*x + 9};
                \addplot[
                    only marks,
                    mark=*,
                    color=oc-red-7
                ]
                coordinates {(1,1) (2,0) (3,1) (4,-1)};

                \node at (axis cs:1,1.5) [anchor=west] {\small $(1,1)$};
                \node at (axis cs:2,0.2) [anchor=south] {\small $(2,0)$};
                \node at (axis cs:3,1.5) [anchor=west] {\small $(3,1)$};
                \node at (axis cs:4,-1) [anchor=west] {\small $(4,-1)$};

                \end{axis}
            \end{tikzpicture}

            \scriptsize\textbf{Illustration:} The Lagrange interpolation polynomial for points $\{(1,1), (2,0), (3,1), (4,-1)\}$ visualized over $\mathbb{R}$.
        \end{center}
    \end{frame}

    \begin{frame}{Recap: QAP}
        \begin{figure}[H]
            \centering
            \begin{tikzpicture}
                \begin{axis}[
                    axis lines = middle,
                    xlabel = {$T$},
                    ylabel = {$y$},
                    ymin = -0.99, ymax = 3.99,
                    xmin = -0.5, xmax = 3.99,
                    domain = 0:5,
                    samples = 100,
                    ytick = {-1,...,4},
                    xtick = {0,1,...,4},
                    zkdlgrid
                ]

                % p(T)
                \addplot[
                    color=oc-blue-7,
                    thick
                ]
                {3/2*x - 1/2*x^2};
                \addplot[
                    only marks,
                    mark=*,
                    color=oc-red-7
                ]
                coordinates {(0, 0) (1, 1) (2, 1) (3, 0)};
                \node at (axis cs:3.35,-0.7) [anchor=east] {$p(T)$};

                \node at (axis cs:0,-0.2) [anchor=west] {\tiny $(0, 0)$};
                \node at (axis cs:1,1) [anchor=south] {\tiny $(1, 1)$};
                \node at (axis cs:2.1,1) [anchor=south] {\tiny $(2, 1)$};
                \node at (axis cs:2.9,0.2) [anchor=west] {\tiny $(3, 0)$};

                % q(T)
                \addplot[
                    color=oc-green-7,
                    thick
                ]
                {1/3*x^3 - 2*x^2 + 8/3*x + 1};
                \addplot[
                    only marks,
                    mark=*,
                    color=oc-red-7
                ]
                coordinates {(0,1) (1,2) (2,1) (3, 0)};
                \node at (axis cs:3.7,0.7) [anchor=south] {$q(T)$};

                \node at (axis cs:0,1) [anchor=west] {\tiny $(0, 1)$};
                \node at (axis cs:1,2) [anchor=south] {\tiny $(1, 2)$};

                % (p+q)(T)
                \addplot[
                    color=oc-orange-7,
                    thick
                ]
                {1/3*x^3 - 5/2*x^2 + 25/6*x + 1};
                \addplot[
                    only marks,
                    mark=*,
                    color=oc-red-7
                ]
                coordinates {(0,1) (1,3) (2,2) (3, 0)};
                \node at (axis cs:3.4,-0.55) [anchor=west] {$(p+q)(T)$};

                \node at (axis cs:1,3) [anchor=south] {\tiny $(1, 3)$};
                \node at (axis cs:2,2) [anchor=west] {\tiny $(2, 2)$};

                \end{axis}
            \end{tikzpicture}
            \caption{Addition of two polynomials}
        \end{figure}
    \end{frame}

    \begin{frame}{Recap: QAP}
        \vspace{-5pt}
        Now, using coefficients encoded with polynomials, we can build a constraint number
        $j \in \{1, \dots, m\}$ in the next way:
        \vspace{-5pt}
        \begin{align*}
            &(z_1\ell_1(j) + z_2\ell_2(j) + \dots + z_n\ell_n(j)) \times \\
            \times &(z_1r_1(j) + z_2r_2(j) + \dots + z_nr_n(j)) = \\
            = &(z_1o_1(j) + z_2o_2(j) + \dots + z_no_n(j))
            \vspace{-5pt}
        \end{align*}
        \pause
        Or written more concisely:
        \vspace{-5pt}
        \begin{equation*}
            \left( \sum_{i = 1}^{n} z_i\ell_i(T) \right) \times \left( \sum_{i = 1}^{n} z_ir_i(T) \right) = \left( \sum_{i = 1}^{n} z_io_i(T) \right)
        \end{equation*}

        \vspace{-5pt}
        \begin{equation*}
            \ell(T) \times r(T) = o(T) \; \commentgray{on $\{1,\dots,m\}$}
        \end{equation*}
    \end{frame}

    \begin{frame}{Recap: QAP}
        Now, we can define a \textbf{master polynomial} $m(T)$, that has zeros at all elements from the set
        $\Omega = \{1,\dots,m\}$
        \vspace{-5pt}
        \begin{equation*}
            m(T) = \ell(T) \times r(T) - o(T)
            \vspace{-5pt}
        \end{equation*}
        \pause
        It means, that $m(T)$ can be divided by \textbf{vanishing polynomial} $z_{\Omega}(T)$ without a remainder!
        \vspace{-8pt}
        \begin{equation*}
            z_{\Omega}(T) = \prod_{i=1}^m (T - i), \qquad \text{$h(T) = \frac{m(T)}{z_{\Omega}(T)}$ is a polynomial}
            \vspace{-5pt}
        \end{equation*}
    \end{frame}

    \section{Encrypted Verification}

    \begin{frame}{Current Point}
        We've managed to encode into \textbf{a single polynomial} an entire computation (a program),
        of any size, independent of how much data it consumes. \\ \pause
        \vspace{10pt}
        Now, we need to figure out the protocol, how a prover can succinctly prove the knowledge of
        a correct witness for some circuit to a verifier, additionally, make it zero-knowledge and
        non-interactive.\pause

        Where the knowledge of the correct witness is a knowledge of the quotient polynomial $h(T)$.
        \begin{equation*}
            m(T) = h(T) \times z_{\Omega}(T)
        \end{equation*}

        \begin{remark}
            Further, for brevity, we will denote $z_{\Omega}(T)$ as $z(T)$.
        \end{remark}
    \end{frame}

    \begin{frame}{Notation Preliminaries: Groups}
        In this section, we will use:
        \begin{itemize}[label=\ding{51}]
            \item Group of points on elliptic curve
            denoted as $\mathbb{G}$ of prime order $q$ with a generator $g$.\pause
            \item The symmetric pairing function $e: \mathbb{G} \times \mathbb{G} \to \mathbb{G}_T$, where
            $(\mathbb{G}_T, \times)$ is a target group (typically, just a scalar from extension $\mathbb{F}_{p^k}$).
        \end{itemize}
        \pause
        \begin{remarkblock}{Recall}
            The core property of the pairing function $e$ is the \textbf{bilinearity}:
            \begin{equation*}
                e(g^{\alpha}, g^{\beta}) = e(g^{\alpha\beta}, g) = e(g, g^{\alpha\beta}) = e(g,g)^{\alpha\beta}.
            \end{equation*}

            Here, $g^{\alpha}$ is the same as ``scalar multiplication of a generator by a scalar $\alpha \in \mathbb{Z}_q$''.
        \end{remarkblock}
    \end{frame}

    \begin{frame}{Notation Preliminaries: QAP}
        Recall that the core equation to be proven is:
        \begin{equation*}
            \textcolor{oc-indigo-8}{\underbrace{\left( \sum_{i = 1}^{n} z_i\ell_i(T) \right)}_{\text{left wires encoding}}} \times \textcolor{oc-yellow-8}{\underbrace{\left( \sum_{i = 1}^{n} z_ir_i(T) \right)}_{\text{right wires encoding}}} - \textcolor{oc-green-8}{\underbrace{\left( \sum_{i = 1}^{n} z_io_i(T) \right)}_{\text{output encodings}}} = z(T)h(T)
        \end{equation*}
        \pause

        Shortly, $\ell(T) \times r(T) - o(T) = z(T)h(T)$.
    \end{frame}

    \begin{frame}{Naive Proof}
        Suppose, we are given a circuit $\Circ$ with a maximum degree $d$ of polynomials
        used underneath.

        Thus, all parties additionally know the target polynomial $z(T)$ and QAP polynomials
        $\{\ell_i(T)\}_{i \in [n]}, \{r_i(T)\}_{i \in [n]}, \{o_i(T)\}_{i \in [n]}$, where $n$ is
        number of witness elements.
        \pause

        \textbf{Prover}
        \vspace{-5pt}
        \begin{itemize}[label=\ding{51}]
            \item \vspace{-3pt} Provides witness $\mathbf{z}$ to a Verifier.
        \end{itemize}
        \pause
        \textbf{Verifier}
        \vspace{-5pt}
        \begin{itemize}[label=\ding{51}]
            \item \vspace{-3pt} Checks $\left( \sum_{i = 1}^{n} z_i\ell_i(T) \right) \times \left( \sum_{i = 1}^{n} z_ir_i(T) \right) = \left( \sum_{i = 1}^{n} z_io_i(T) \right)$ on $\Omega$.
        \end{itemize}
    \end{frame}

    \begin{frame}{Naive Proof}
        \begin{itemize}
            \item[\xmark] Succinct
            \item[\cmark] Non-Interactive
            \item[\xmark] Zero-Knowledge
        \end{itemize}
        \pause

        \vspace{20pt}
        \begin{columns}
            \begin{column}{0.7\textwidth}
                The verifier could actually just run a program that represents a circuit
                $\Circ$ on witness data $\mathbf{z}$.
            \end{column}

            \begin{column}{0.2\textwidth}
                \includegraphics[width=\linewidth]{images/common/new-moon-with-face.jpg}
            \end{column}
        \end{columns}\pause

        We, definitely, need to encrypt the witness data $\mathbf{z}$ somehow...
    \end{frame}

    \begin{frame}{Encryption and KZG Commitment}
        Let us define the \textit{encryption} operation as follows:
        \vspace{-10pt}
        \begin{equation*}
            \mathsf{Enc}: \mathbb{Z}_q \to \mathbb{G}, \quad \mathsf{Enc}(x) := g^x
            \vspace{-5pt}
        \end{equation*}\pause
        Essentially, $\mathsf{Enc}(p(\tau))$ is the \textbf{KZG Commitment} $\mathsf{com}(p)$.
        \begin{example}
            Consider the polynomial: $p(T) = T^2 - 5T + 2$, the encryption of $p(\tau)$ can be found as follows:
            {\large
            \begin{equation*}
                \mathsf{Enc}(p(\tau)) = g^{p(\tau)} = g^{\left( \tau^2 - 5\tau + 2 \right)}  = \left( g^{\tau^2} \right)^1 \cdot {\left( g^{\tau^1} \right)}^{-5} \cdot \left(g^{\tau^0}\right)^2
            \end{equation*}}
        \end{example}\pause
        \begin{questionblock}{Question}
            KZG Commitment requires encrypted powers of $\tau$: $\{g^{\tau^i}\}_{i \in [d]}$. But
            where can the prover take them?
        \end{questionblock}
    \end{frame}

    \begin{frame}{Trusted Setup}
        \textbf{Trusted Party Setup}
        \vspace{-5pt}
        \begin{itemize}[label=\ding{51}]
            \item \vspace{-3pt} Picks a random value $\tau \leftarrowS \mathbb{F}$. \pause
            \item \vspace{-3pt} Calculates the public parameters $\{g^{\tau^i}\}_{i \in [d]}$.\pause
            \item \vspace{-3pt} \textbf{Deletes} $\tau$ (toxic waste).\pause
            \item \vspace{-3pt} \textbf{Outputs} prover parameters $\mathsf{pp} \gets \{g^{\tau^i}\}_{i \in [d]}$ and verifier parameters $\mathsf{vp} \gets \mathsf{com}(z)$.
        \end{itemize} \pause

        This way, we can find the KZG commitment for each polynomial. For example:
        \begin{equation*}
            \mathsf{com}(\ell) \triangleq g^{\ell(\tau)} = g^{\sum_{i=0}^d \ell_i \tau^i} = \prod_{i=0}^d (g^{\tau^i})^{\ell_i},
        \end{equation*}
    \end{frame}

    \begin{frame}{Encrypted Verification}
        Now, we can calculate the following KZG commitments (or, synonymously, encryptions):
        \begin{equation*}
            g^{\ell(\tau)}, g^{r(\tau)}, g^{o(\tau)}, g^{h(\tau)}, g^{z(\tau)}
        \end{equation*}
        But how can we verify $h(T)z(T) = \ell(T)r(T) - o(T)$ in the encrypted space?
        \pause

        Well, first notice that, according to the \textbf{Schwartz-Zippel Lemma}, \textit{with overwhelming probability} the check is equivalent to:
        \begin{equation*}
            \ell(\tau)r(\tau) = z(\tau)h(\tau) + o(\tau).
        \end{equation*}
        \pause
        So, we can check this equality as follows:
        \begin{equation*}
            e(\mathsf{com}(\ell), \mathsf{com}(r)) = e(\mathsf{com}(z), \mathsf{com}(h)) \cdot e(\mathsf{com}(o), g),
        \end{equation*}
    \end{frame}

    \begin{frame}{Attempt \#1: Encrypted Verification}
        \scriptsize

        \textbf{Trusted Setup: }
        $\tau \leftarrowS \mathbb{F} \textbf{,} \quad \{g^{\tau^i}\}_{i \in [d]}\textbf{,} \quad g^{z(\tau)} \textbf{,} \quad \textbf{delete}(\tau)$.
        \vspace{10pt}

        \begin{center}
            \begin{tikzpicture}
                \node[inner sep=0pt, align=center] (prover) at (-3.5, -3) {
                    \includegraphics[width=2cm]{images/common/prover.png}\\Prover $\mathcal{P}$
                };

                \node[inner sep=0pt, align=center] (verifier) at (4, -3) {
                    \includegraphics[width=2cm]{images/common/verifier.png}\\Verifier $\mathcal{V}$
                };

                \onslide<2->{
                    \node[inner sep=0pt, align=center] (pdata) at (-3, -0.6) {
                        \begin{minipage}{4cm}
                            \begin{itemize}[label=\ding{51}]
                                \item $h(T) = \frac{\ell(T) \times r(T) - o(T)}{z(T)}$.
                                \onslide<3->{\item {
                                    KZG commitments:
                                    \vspace{-10pt}
                                    \begin{align*}
                                        \pi_L &\gets \mathsf{com}(\ell), &\pi_R &\gets \mathsf{com}(r), \\
                                        \pi_O &\gets \mathsf{com}(o), &\pi_H &\gets \mathsf{com}(h),
                                    \end{align*}
                                }}
                            \end{itemize}
                        \end{minipage}
                    };
                }

                \onslide<4->{
                    \draw[arrow] (prover) -- (verifier) node[midway, above=1mm] {$\boldsymbol{\pi} = (\pi_L,\pi_R,\pi_O,\pi_H)$};
                }

                \onslide<5->{
                    \node[inner sep=0pt, align=center] (vdata) at (5, -1.1) {
                        \begin{minipage}{6cm}
                            \begin{itemize}[label=\ding{51}]
                                \item $e(\pi_L, \pi_R) \xlongequal{?} \\ e(\mathsf{com}(z), \pi_H) \cdot e(\pi_O, g)$.
                            \end{itemize}
                        \end{minipage}
                    };
                }
            \end{tikzpicture}
        \end{center}
    \end{frame}

    \begin{frame}{Attempt \#1: Encrypted Verification}
        \begin{itemize}
            \item[\cmark] Succinct
            \item[\cmark] Non-Interactive
            \item[\cmark] Zero-Knowledge \pause
            \item[\xmark] \vspace{7pt} Does it work?
        \end{itemize}
    \end{frame}

    \begin{frame}{Why it doesn't work??}\begin{adjustwidth}{-5mm}{0mm}
        \scriptsize

        \textbf{Trusted Setup: }
        $\tau \leftarrowS \mathbb{F} \textbf{,} \quad \{g^{\tau^i}\}_{i \in [d]}\textbf{,} \quad g^{z(\tau)}\textbf{,}$ \textbf{delete}($\tau$).
        \vspace{10pt}

        \begin{center}
            \begin{tikzpicture}
                \node[inner sep=0pt, align=center] (prover) at (-3.5, -2.4) {
                    \includegraphics[width=2cm]{images/common/prover.png}\\Prover $\mathcal{P}$
                };

                \node[inner sep=0pt, align=center] (verifier) at (4, -2.4) {
                    \includegraphics[width=2cm]{images/common/verifier.png}\\Verifier $\mathcal{V}$
                };

                \onslide<2->{\node[inner sep=0pt, align=center] (pdata) at (-2.3, 0) {
                    \begin{minipage}{6cm}
                        \begin{itemize}[label=\ding{51}]
                            \onslide<3->{\item $h'(T) \leftarrowS \mathbb{F}[T], \quad m'(T) = z(T) \times h'(T)$.}
                            \onslide<4->{\item Finds $\ell'(T), r'(T), o'(T)$ such that: \\ $\ell'(T) \times r'(T) - o'(T) = m'(T)$}
                            \onslide<5->{\item {
                                KZG commitments:
                                \vspace{-10pt}
                                \begin{align*}
                                    \pi_{\ell'} &\gets \mathsf{com}(\ell'), &\pi_{r'} &\gets \mathsf{com}(r'), \\
                                    \pi_{o'} &\gets \mathsf{com}(o'), &\pi_{h'} &\gets \mathsf{com}(h'),
                                \end{align*}
                            }}
                        \end{itemize}
                    \end{minipage}
                };}

                \onslide<6->{\draw[arrow] (prover) -- (verifier) node[midway, above=1mm] {$\boldsymbol{\pi} = (\pi_{\ell'},\pi_{r'},\pi_{o'},\pi_{h'})$};}

                \onslide<7->{\node[inner sep=0pt, align=center] (vdata) at (5.1, -0.6) {
                    \begin{minipage}{6cm}
                        \begin{itemize}[label=\ding{51}]
                            \item $e(\pi_{\ell'}, \pi_{r'}) \xlongequal{?} \\ e(\mathsf{com}(z), \pi_{h'}) \cdot e(\pi_{o'}, g)$.
                        \end{itemize}
                    \end{minipage}
                };}
            \end{tikzpicture}
        \end{center}

        \pause
        \begin{alertblock}{Problem}
            Prover isn't forced to use the values from the trusted setup.
        \end{alertblock}
    \end{adjustwidth}\end{frame}

    \begin{frame}{Attempt \#2: Proof Of Exponent}
        \scriptsize

        \textbf{Trusted Setup: }
        $\tau, \chg{\alpha} \leftarrowS \mathbb{F} \textbf{,} \quad \{ \{g^{\tau^i}, \chg{g^{\alpha \tau^i}}\}_{i \in [d]} \} \textbf{,} \quad \{g^{z(\tau)}, \chg{g^{\alpha}}\}\textbf{,} \quad \textbf{delete}(\tau, \chg{\alpha})$.
        \vspace{10pt}

        \begin{center}
            \begin{tikzpicture}
                \node[inner sep=0pt, align=center] (prover) at (-3.5, -3) {
                    \includegraphics[width=2cm]{images/common/prover.png}\\Prover $\mathcal{P}$
                };

                \node[inner sep=0pt, align=center] (verifier) at (4, -3) {
                    \includegraphics[width=2cm]{images/common/verifier.png}\\Verifier $\mathcal{V}$
                };

                \onslide<2->{\node[inner sep=0pt, align=center] (pdata) at (-3, -0.6) {
                    \begin{minipage}{4cm}
                        {\tiny\begin{itemize}[label=\ding{51}]
                            \item $h(T) = \frac{\ell(T) \times r(T) - o(T)}{z(T)}$.
                            \onslide<3->{\item {
                                KZG commitments:
                                \vspace{-8pt}
                                \begin{align*}
                                    \pi_L &\gets g^{\ell(\tau)}, \quad & \chg{\pi_L'} &\gets \chg{g^{\alpha \ell(\tau)}},\\[-2pt]
                                    \pi_R &\gets g^{r(\tau)}, \quad & \chg{\pi_R'} &\gets \chg{g^{\alpha r(\tau)}},\\[-2pt]
                                    \pi_O &\gets g^{o(\tau)}, \quad & \chg{\pi_O'} &\gets \chg{g^{\alpha o(\tau)}},\\[-2pt]
                                    \pi_H &\gets g^{h(\tau)}, \quad & \chg{\pi_H'} &\gets \chg{g^{\alpha h(\tau)}}.
                                \end{align*}
                            }}
                        \end{itemize}}
                    \end{minipage}
                };}

                \onslide<4->{\draw[arrow] (prover) -- (verifier) node[midway, above=1mm] {$\boldsymbol{\pi} = (\pi_L,\pi_R,\pi_O,\pi_H,\chg{\pi_L',\pi_R',\pi_O',\pi_H'})$};}

                \onslide<5->{\node[inner sep=0pt, align=center] (vdata) at (5, -0.6) {
                    \begin{minipage}{4cm}
                        {\tiny\begin{itemize}[label=\ding{51}]
                            \item $e(\pi_L, \pi_R) \xlongequal{?}\\ e(\mathsf{com}(z), \pi_H) \cdot e(\pi_O, g)$.
                            \onslide<6->{\item \chg{
                                Proof of Exponent:
                                \vspace{-8pt}
                                \begin{align*}
                                    e(\pi_L, g^{\alpha}) &= e(\pi_L', g), \qquad\qquad\qquad\\[-2pt]
                                    e(\pi_R, g^{\alpha}) &= e(\pi_R', g), \\[-2pt]
                                    e(\pi_O, g^{\alpha}) &= e(\pi_O', g), \\[-2pt]
                                    e(\pi_H, g^{\alpha}) &= e(\pi_H', g).
                                \end{align*}
                            }}
                        \end{itemize}}
                    \end{minipage}
                };}
            \end{tikzpicture}
        \end{center}
    \end{frame}

    \begin{frame}{Including PoE}
        \begin{itemize}
            \item[\cmark] Succinct
            \item[\cmark] Non-Interactive
            \item[\cmark] Zero-Knowledge \pause
            \item[\xmark] Sound
        \end{itemize}
        \pause

        \vspace{15pt}
        \begin{alertblock}{Problem}
            There is no guarantee that the same witness $\mathbf{z}$ was used to calculate all the commitments $\pi_L, \pi_R, \pi_O, \pi_H$.
        \end{alertblock}
    \end{frame}

    \section{Make It Sound}

    \begin{frame}{Additional Optimization}
        Recall that:
        \begin{equation*}
            \ell(T) = \sum_{i=1}^n z_i\ell_i(T), \quad r(T) = \sum_{i=1}^n z_i r_i(T), \quad o(T) = \sum_{i=1}^n z_io_i(T).
        \end{equation*}
        \pause

        Here public data is:
        \begin{equation*}
            \{\ell_i(T)\}_{i \in [n]}, \{r_i(T)\}_{i \in [n]}, \{o_i(T)\}_{i \in [n]}
        \end{equation*}
        \pause

        Moreover, it's defined only by the circuit and trusted setup, thus, it can be calculated before proof generation as a part of the trusted setup.
    \end{frame}

    \begin{frame}{Additional Optimization}
        Updated Trusted Setup:
        \vspace{-10pt}
        {\scriptsize \begin{align*}
            \{g^{\tau^i}\}_{i \in [d]}&, \quad \{g^{\alpha \tau^i}\}_{i \in [d]}, \\
            \chg{\{g^{\ell_i(\tau)}\}_{i \in [n]}}&, \quad \chg{\{g^{\alpha \ell_i(\tau)}\}_{i \in [n]}},\\
            \chg{\{g^{r_i(\tau)}\}_{i \in [n]}}&, \quad \chg{\{g^{\alpha r_i(\tau)}\}_{i \in [n]}},\\
            \chg{\{g^{o_i(\tau)}\}_{i \in [n]}}&, \quad \chg{\{g^{\alpha o_i(\tau)}\}_{i \in [n]}}
        \end{align*}}
        \pause

        Consider the polynomial $\ell(T) = \sum_{i=1}^n z_i\ell_i(T)$.
        \pause

        $\mathcal{P}$ can compute the KZG commitment $\pi_L$ and its PoE $\pi_L'$ as follows:
        \begin{align*}
            \pi_L &\triangleq g^{\ell(\tau)} = g^{\sum_{i=1}^n z_i \ell_i(\tau)} = \prod_{i=1}^n (g^{\ell_i(\tau)})^{z_i}, \\
            \pi_L' &\triangleq g^{\alpha \ell(\tau)} = g^{\alpha \sum_{i=1}^n z_i \ell_i(\tau)} = \prod_{i=1}^n (g^{\alpha \ell_i(\tau)})^{z_i}.
        \end{align*}
    \end{frame}

    \begin{frame}{Witness Consistency Check}
        To prove that the same $\mathbf{z}$ is used in all commitments, we need some ``checksum''
        term that will somehow combine all polynomials $\ell(T)$, $r(T)$, and $o(T)$ with the witness
        $\mathbf{z}$.
        \pause

        We can do so by proving the next simple statement:
        \vspace{-5pt}
        \begin{equation*}
            g^{\ell(\tau) + r(\tau) + o(\tau)} = \pause \prod_{i=1}^n {\left(g^{\ell_i(\tau) + r_i(\tau) + o_i(\tau)}\right)}^{z_i}
            \vspace{-5pt}
        \end{equation*}
        \pause
        And we already know how to do that --- PoE!
    \end{frame}

    \begin{frame}{Witness Consistency Check}
        Let us introduce one more coefficient...
        \vspace{-5pt}
        \begin{equation*}
            \beta \leftarrowS \mathbb{F}
            \vspace{-5pt}
        \end{equation*}
        \pause
        Extended trusted setup contains additional values:
        \vspace{-5pt}
        \begin{equation*}
            g^{\beta}, \quad \{g^{\beta(\ell_i(\tau) + r_i(\tau) + o_i(\tau))}\}_{i \in [n]}
            \vspace{-5pt}
        \end{equation*}
        \pause
        Prover needs to calculate $\pi_{\beta}$:
        \vspace{-5pt}
        \begin{equation*}
            \pi_{\beta} \gets g^{\beta(\ell(\tau) + r(\tau) + o(\tau))} = \pause \prod_{i=1}^{n} \left(g^{\beta(\ell_i(\tau) + r_i(\tau) + o_i(\tau))}\right)^{z_i}
        \end{equation*}
        \pause
        And easy check for verifier:
        \vspace{-5pt}
        \begin{equation*}
            e(\pi_L\pi_R\pi_O, g^{\beta}) = e(\pi_{\beta}, g).
        \end{equation*}
    \end{frame}

    \begin{frame}{One more time... that doesn't work}
        If the witness is consistent, the following condition must hold:
        \vspace{-3pt}
        {\small \begin{equation*}
            (z_{\ell,i}\ell_i(\tau) + z_{r,i}r_i(\tau) + z_{o,i}o_i(\tau))\beta = z_{\beta,i}\beta(\ell_i(\tau) + r_i(\tau) + o_i(\tau)) \; \commentgray{for all $i \in [n]$}
            \vspace{-10pt}
        \end{equation*}}
        \pause

        But, what if $\ell_i \equiv r_i$. Let us call them $q$, thus:
        \vspace{-3pt}
        \begin{equation*}
            (z_{\ell,i} + z_{r,i})q + z_{o,i}o_i(\tau)= z_{\beta,i}(2q + o_i(\tau)) \; \commentgray{for all $i \in [n]$}
            \vspace{-10pt}
        \end{equation*}
        \pause

        The adversary can choose $z_{\ell,i}$, $z_{r,i}$ and $z_{o,i}$ such that:
        \vspace{-5pt}
        \begin{equation*}
            z_{\beta,i} := z_{o,i} \quad \text{and} \quad z_{\ell,i}=2z_{o,i} - z_{r,i}
            \vspace{-5pt}
        \end{equation*}
        \pause

        \begin{example}
            \begin{equation*}
                z_{\beta} = z_o = 5, \quad z_{\ell} = 7, \quad z_r = 3
                \vspace{-25pt}
            \end{equation*}

            \begin{align*}
                (7 + 3)q + 5o(\tau) = 5(2q + o(\tau)) \\
                10q + 5o(\tau) = 10q + 5o(\tau)
            \end{align*}
        \end{example}
    \end{frame}

    \begin{frame}{More coefficients!}
        The main problem is that if $r_i \equiv \ell_i$ then $\beta r_i \equiv \beta \ell_i$.\pause

        Let us fix that by introducing separate $\beta$ coefficients for $\ell$, $r$ and $o$.
        \begin{equation*}
            \left( \beta_L, \beta_R, \beta_O \right) \leftarrowS \mathbb{F}^3
            \vspace{-5pt}
        \end{equation*}
        \pause

        Therefore, $\beta_R r_i \neq \beta_L \ell_i$ even if $r_i \equiv \ell_i$, so the previous hack doesn't work.\pause

        So, finally, the trusted setup is updated with:
        \vspace{-5pt}
        \begin{equation*}
            g^{\beta_L}, g^{\beta_R}, g^{\beta_O}, \{g^{\beta_L\ell_i(\tau) + \beta_Rr_i(\tau) + \beta_Oo_i(\tau)}\}_{i \in [n]}
            \vspace{-5pt}
        \end{equation*}
        \pause

        Verification:
        \vspace{-5pt}
        \begin{equation*}
            e(\pi_L, g^{\beta_L}) \cdot e(\pi_R, g^{\beta_R}) \cdot e(\pi_O, g^{\beta_O}) = e(\pi_{\beta}, g)
        \end{equation*}
    \end{frame}

    \begin{frame}{Even more coefficients!}
        As the adversary has access to the public $g^{\beta_L}$, $g^{\beta_R}$, $g^{\beta_O}$ he still can cheat the verifier by calculating a modified $\pi_{\beta}$.

        \begin{example}
            \small
            Consider a constraint $z_1 \times z_1 = z_2$. Let us try to assign $2$ and $5$ for $z_1$ in a single constraint. As $2 \times 5 = 10$, $z_2$ should contain the value $10$.
            \vspace{-5pt}
            \begin{equation*}
                \mathbf{z} = (z_1, z_2) = (2, 10)
                \vspace{-5pt}
            \end{equation*}
            \pause
            The next QAP can be built:
            \vspace{-5pt}
            \begin{align*}
                \ell(T) &= 2\ell_1(T) + 10\ell_2(T) \\
                r(T) &= 2r_1(T) + 3 + 10r_2(T) \\
                o(T) &= 2o_1(T) + 10o_2(T)
            \end{align*}
            \pause
            Compute $\pi_{\beta}$ as:
            \vspace{-5pt}
            \begin{equation*}
                (g^{\beta_L\ell_1(\tau) + \beta_Rr_1(\tau) + \beta_Oo_1(\tau)})^2 \cdot (g^{\beta_R})^3 \cdot (g^{\beta_L\ell_2(\tau) + \beta_Rr_2(\tau) + \beta_Oo_2(\tau)})^{10}
            \end{equation*}
        \end{example}
    \end{frame}

    \begin{frame}{Even more coefficients!}
        To prevent this, let us introduce... one more coefficient!
        \vspace{-5pt}

        \begin{equation*}
            \gamma \leftarrowS \mathbb{F}
            \vspace{-5pt}
        \end{equation*}
        \pause

        So, finally... the trusted setup is updated with:
        \vspace{-5pt}
        \begin{equation*}
            g^{\gamma}, g^{\beta_L\gamma}, g^{\beta_R\gamma}, g^{\beta_O\gamma}, \{g^{\beta_L\ell_i(\tau) + \beta_Rr_i(\tau) + \beta_Oo_i(\tau)}\}_{i \in [n]}
            \vspace{-5pt}
        \end{equation*}
        \pause

        Proving process isn't changed, unlike verification:
        \begin{equation*}
            e(\pi_L, g^{\beta_L\gamma}) \cdot e(\pi_R, g^{\beta_R\gamma}) \cdot e(\pi_O, g^{\beta_O\gamma}) = e(\pi_{\beta}, g^{\gamma})
        \end{equation*}

        That makes it infeasible to cheat.
    \end{frame}

    \begin{frame}{Attempt \#3: Sound SNARK Protocol}
        \scriptsize

        \textbf{Trusted Setup: }
        $\tau, \alpha, \chg{\beta_L, \beta_R, \beta_O, \gamma} \leftarrowS \mathbb{F} \textbf{,} \quad \{ \{g^{\tau^i}, g^{\alpha \tau^i} \}_{i \in [d]}\textbf{,} \quad \chg{\{g^{\beta_L\ell_i(\tau) + \beta_Rr_i(\tau) + \beta_Oo_i(\tau)}\}_{i \in [n]}} \} \textbf{,}$
        $\{ g^{z(\tau)}, g^{\alpha}, \chg{g^{\beta_L}, g^{\beta_R}, g^{\beta_O}, g^{\beta_L\gamma}, g^{\beta_R\gamma}, g^{\beta_O\gamma}, g^{\gamma}}\}\textbf{,} \quad \textbf{delete}(\tau, \alpha, \chg{\beta_L, \beta_R, \beta_O, \gamma})$.
        \vspace{10pt}

        \begin{center}
            \begin{tikzpicture}
                \node[inner sep=0pt, align=center] (prover) at (-3.5, -3) {
                    \includegraphics[width=2cm]{images/common/prover.png}\\Prover $\mathcal{P}$
                };

                \node[inner sep=0pt, align=center] (verifier) at (4, -3) {
                    \includegraphics[width=2cm]{images/common/verifier.png}\\Verifier $\mathcal{V}$
                };

                \onslide<2->{\node[inner sep=0pt, align=center] (pdata) at (-3, 0) {
                    \begin{minipage}{4cm}
                        {\tiny\begin{itemize}[label=\ding{51}]
                            \item $h(T) = \frac{\ell(T) \times r(T) - o(T)}{z(T)}$.
                            \onslide<3->{\item {
                                KZG commitments:
                                \vspace{-8pt}
                                \begin{align*}
                                    \pi_L &\gets g^{\ell(\tau)}, \quad &\pi_L' &\gets g^{\alpha \ell(\tau)},\\[-2pt]
                                    \pi_R &\gets g^{r(\tau)}, \quad &\pi_R' &\gets g^{\alpha r(\tau)},\\[-2pt]
                                    \pi_O &\gets g^{o(\tau)}, \quad &\pi_O' &\gets g^{\alpha o(\tau)},\\[-2pt]
                                    \pi_H &\gets g^{h(\tau)}, \quad &\pi_H' &\gets g^{\alpha h(\tau)}.
                                \end{align*}
                                \vspace{-16pt}
                            }}
                            \onslide<4->{\item \chg{
                                \vspace{-18pt}
                                $\pi_{\beta} \gets g^{\beta_L\ell(\tau) + \beta_Rr(\tau) + \beta_Oo(\tau)}$
                            }}
                        \end{itemize}}
                    \end{minipage}
                };}

                \onslide<5->{\draw[arrow] (prover) -- (verifier) node[midway, above=1mm] {$\boldsymbol{\pi} = (\pi_L,\pi_R,\pi_O,\pi_H,\pi_L',\pi_R',\pi_O',\pi_H',\chg{\pi_{\beta}})$};}

                \onslide<6->{\node[inner sep=0pt, align=center] (vdata) at (5, -0.2) {
                    \begin{minipage}{5cm}
                        {\tiny\begin{itemize}[label=\ding{51}]
                            \item $e(\pi_L, \pi_R) \xlongequal{?}\\ e(\mathsf{com}(z), \pi_H) \cdot e(\pi_O, g)$.
                            \onslide<7->{\item {
                                Proof of Exponent:
                                \vspace{-8pt}
                                \begin{align*}
                                    e(\pi_L, g^{\alpha}) &= e(\pi_L', g), \qquad\qquad\qquad\\[-2pt]
                                    e(\pi_R, g^{\alpha}) &= e(\pi_R', g), \\[-2pt]
                                    e(\pi_O, g^{\alpha}) &= e(\pi_O', g), \\[-2pt]
                                    e(\pi_H, g^{\alpha}) &= e(\pi_H', g).
                                \end{align*}
                                \vspace{-16pt}
                            }}
                            \onslide<8->{\item \chg{
                                $e(\pi_L, g^{\gamma\beta_L}) \cdot e(\pi_R, g^{\gamma\beta_R}) \cdot \\ \cdot e(\pi_O, g^{\gamma\beta_O}) = e(\pi_{\beta}, g^{\gamma})$
                            }}
                        \end{itemize}}
                    \end{minipage}
                };}
            \end{tikzpicture}
        \end{center}
    \end{frame}

    \section{Make it Zero-Knowledge}

    \begin{frame}{Why do we need modifications?}
        \begin{questionblock}{Question}
            \begin{itemize}[label=\ding{228}]
                \item According to the security assumptions, given any element of our proof $\boldsymbol{\pi}$, it is impossible to restore the witness $\mathbf{z}$.\pause
                \item Well, that is true, but does it guarantee that no other information about $\mathbf{z}$ is \textbf{leaked}?\pause
            \end{itemize}

            It does not! For example, given $\boldsymbol{\pi}$, I can check if $\ell(T) = 10r(T)$. \pause

            How? Simply check
            \begin{equation*}
                \pi_L \xlongequal{?} \pi_R^{10}.
            \end{equation*}
            \pause

            This works since $\pi_L = g^{\ell(\tau)}$ and $\pi_R = g^{r(\tau)}$, so
            \begin{equation*}
                \pi_L = \pi_R^{10} \Leftrightarrow g^{\ell(\tau)} = g^{10r(\tau)} \Leftrightarrow \ell(\tau) = 10r(\tau) \Leftrightarrow \ell(T) = 10r(T)
            \end{equation*}
        \end{questionblock}
    \end{frame}

    \begin{frame}{How do we fix that? Analysis}
        \begin{block}{Analysis}
            \begin{itemize}[label=\ding{228}]
                \item Currently, proof consists of four ingredients: $\pi_L, \pi_R, \pi_O, \pi_H$, and their corresponding PoEs. \pause
                \item But, current construction creates some ``connection'' between these ingredients. For example, $\pi_L^{101} = \pi_R$ reveals whether $101\ell(T)=r(T)$ or $\pi_L\pi_O=\pi_R$ whether $\ell(T) + o(T) = r(T)$. \pause
                \item Therefore, the prover must add some ``noise'' to the proof, so that the verifier can't use the proof to extract any information about the witness. The randomness of prover is kept secret. \pause
                \item Simultaneously, this noise would still make the same proof checkable by already defined verification equations.
            \end{itemize}
        \end{block}
        \pause

        \begin{questionblock}{Question}
            So how do we represent such ``noise'' while preserving ``homomorphic'' properties of the proof?
        \end{questionblock}
    \end{frame}

    \begin{frame}{How do we fix that? Doing stuff}
        \begin{block}{Idea \#1}
            Let prover pick random values $\delta_L, \delta_R, \delta_O, \delta_H \leftarrowS \mathbb{F}$ and calculate the ``distorted'' values:
            \begin{equation*}
                \pi_L \gets g^{\ell(\tau) + \delta_L}, \quad \pi_R \gets g^{r(\tau) + \delta_R}, \quad \text{same for $\pi_O, \pi_H$}
            \end{equation*}
            \pause\textbf{Problem (informally):} To make previous verification mechanism work, $\delta_H$ must be picked ``carefully''. However, it's impossible to do so with this construction.
        \end{block}
        \pause

        \begin{block}{Idea \#2}
            Let prover pick random values $\delta_L, \delta_R, \delta_O \leftarrowS \mathbb{F}$ and calculate:
            \begin{equation*}
                \pi_L \gets g^{\ell(\tau) + \delta_L z(\tau)}, \quad \pi_R \gets g^{r(\tau) + \delta_R z(\tau)}, \quad \pi_O \gets g^{o(\tau) + \delta_O z(\tau)}
            \end{equation*}
        \end{block}
    \end{frame}

    \begin{frame}{Making Idea Practical}
        This expression can be easily evaluated. For example,
        \begin{equation*}
            \pi_L = g^{\ell(\tau)} \cdot \left(g^{z(\tau)}\right)^{\delta_L} = \mathsf{com}(\ell) \cdot \mathsf{com}(z)^{\delta_L}
        \end{equation*}
        \pause

        \begin{alertblock}{Problem}
            However, the verifier can't check the equation $e(\pi_L, \pi_R) = e(\mathsf{com}(z), \pi_H) \cdot e(\pi_O, g)$ anymore.
        \end{alertblock}
        \pause

        \begin{block}{Idea \#3}
            Our ``noise'' must be controlled! This can be achieved by calculating $\pi_H$ as $g^{h(\tau) + \Delta_H}$ where $\Delta_H$ depends on $\delta_L, \delta_R, \delta_O$ (and possibly other public parameters).
        \end{block}
    \end{frame}

    \begin{frame}{Making Idea Practical}
        Check $e(\pi_L, \pi_R) = e(\mathsf{com}(z), \pi_H) \cdot e(\pi_O, g)$ is now equivalent to:
        \begin{equation*}
            (\ell(T) + \delta_L z(T))(r(T) + \delta_R z(T)) = (h(T) + \Delta_H)z(T) + (o(T) + \delta_O z(T)),
        \end{equation*}
        \pause

        which, by expanding, gives us the following equation:
        \begin{align*}
            \cancel{\ell(T)r(T)} + \delta_R \ell(T)z(T) + \delta_L z(T)r(T) + \delta_L \delta_R z(T)^2 \\= \cancel{h(T)z(T) + o(T)} + \Delta_Hz(T) + \delta_O z(T)
        \end{align*}
        \pause

        where we can cancel out $\ell(T)r(T)$ and $h(T)z(T) + o(T)$ terms since they are equal based on initial construction. This way, we get the following expression for $\Delta_H$:
        \begin{equation*}
            \boxed{\Delta_H = \delta_O + \delta_R \ell(T) + \delta_L r(T) + \delta_L \delta_R z(T)}
        \end{equation*}
    \end{frame}

    \begin{frame}{Witness Consistency Proof}
        Finally, let us not forget about $\pi_{\beta}$! Previously, we had:
        \begin{equation*}
            \pi_{\beta} = g^{\beta_L\ell(\tau) + \beta_Rr(\tau) + \beta_Oo(\tau)}
        \end{equation*}
        \pause

        Now, we change:
        \begin{itemize}[label=\ding{228}]
            \item $\ell(\tau) \mapsto \ell(\tau) + \delta_L z(\tau)$
            \item $r(\tau) \mapsto r(\tau) + \delta_R z(\tau)$
            \item $o(\tau) \mapsto o(\tau) + \delta_O z(\tau)$
        \end{itemize}
        \pause

        Therefore, our new $\pi_{\beta}$ becomes:
        \begin{equation*}
            \pi_{\beta} = \chg{\left(g^{\beta_Lz(\tau)}\right)^{\delta_L}\left(g^{\beta_R z(\tau)}\right)^{\delta_R}\left(g^{\beta_O z(\tau)}\right)^{\delta_O}}g^{\beta_L \ell(\tau) + \beta_R r(\tau) + \beta_O o(\tau)}
        \end{equation*}
    \end{frame}

    \begin{frame}{Overall Protocol}
        \scriptsize

        \textbf{Trusted Setup: }
        $\tau, \alpha, \beta_L, \beta_R, \beta_O, \gamma \leftarrowS \mathbb{F} \textbf{,} \quad \{ \{g^{\tau^i}, g^{\alpha \tau^i} \}_{i \in [d]}\textbf{,} \quad \{\chg{g^{\beta_L\ell_i(\tau)}, g^{\beta_Rr_i(\tau)}, g^{\beta_Oo_i(\tau)}}\}_{i \in [n]} \} \textbf{,}$
        $\{ g^{z(\tau)}, g^{\alpha}, g^{\beta_L}, g^{\beta_R}, g^{\beta_O}, g^{\beta_L\gamma}, g^{\beta_R\gamma}, g^{\beta_O\gamma}, g^{\gamma}\}\textbf{,} \quad \textbf{delete}(\tau, \alpha, \beta_L, \beta_R, \beta_O, \gamma)$.
        \vspace{10pt}

        \begin{center}
            \begin{tikzpicture}
                \node[inner sep=0pt, align=center] (prover) at (-3.5, -3) {
                    \includegraphics[width=2cm]{images/common/prover.png}\\Prover $\mathcal{P}$
                };

                \node[inner sep=0pt, align=center] (verifier) at (4, -3) {
                    \includegraphics[width=2cm]{images/common/verifier.png}\\Verifier $\mathcal{V}$
                };

                \onslide<2->{\node[inner sep=0pt, align=center] (pdata) at (-3, 0) {
                    \begin{minipage}{4cm}
                        {\tiny\begin{itemize}[label=\ding{51}]
                            \item $h(T) = \frac{\ell(T) \times r(T) - o(T)}{z(T)}$.
                            \onslide<3->{\item {
                                Sample $\delta_L,\delta_R,\delta_O \leftarrowS \mathbb{F}$, compute:
                                \vspace{-8pt}
                                \begin{align*}
                                    \pi_L \gets g^{\ell(\tau)}\chg{(g^{z(\tau)})^{\delta_L}}, \pi_L' \gets g^{\alpha \ell(\tau)}\chg{(g^{\alpha z(\tau)})^{\delta_L}},\\[-2pt]
                                    \pi_R \gets g^{r(\tau)}\chg{(g^{z(\tau)})^{\delta_R}}, \pi_R' \gets g^{\alpha r(\tau)}\chg{(g^{\alpha z(\tau)})^{\delta_R}},\\[-2pt]
                                    \pi_O \gets g^{o(\tau)}\chg{(g^{z(\tau)})^{\delta_O}}, \pi_O' \gets g^{\alpha o(\tau)}\chg{(g^{\alpha z(\tau)})^{\delta_O}},\\[-2pt]
                                    \pi_H \gets g^{h(\tau)}\chg{(g^{\delta_O})(g^{r(\tau)})^{\delta_L}(g^{\ell(\tau)})^{\delta_R}(g^{z(\tau)})^{\delta_L\delta_R}}
                                \end{align*}
                                \vspace{-16pt}
                            }}
                            \onslide<4->{
                                \vspace{-5pt}
                                $\pi_{\beta} = \dots$
                            }
                        \end{itemize}}
                    \end{minipage}
                };}

                \onslide<5->{\draw[arrow] (prover) -- (verifier) node[midway, above=1mm] {$\boldsymbol{\pi} = (\pi_L,\pi_R,\pi_O,\pi_H,\pi_L',\pi_R',\pi_O',\pi_H',\pi_{\beta})$};}

                \onslide<6->{\node[inner sep=0pt, align=center] (vdata) at (5, -0.2) {
                    \begin{minipage}{5cm}
                        {\tiny\begin{itemize}[label=\ding{51}]
                            \item $e(\pi_L, \pi_R) \xlongequal{?}\\ e(\mathsf{com}(z), \pi_H) \cdot e(\pi_O, g)$.
                            \onslide<7->{\item {
                                Proof of Exponent:
                                \vspace{-8pt}
                                \begin{align*}
                                    e(\pi_L, g^{\alpha}) &= e(\pi_L', g), \qquad\qquad\qquad\\[-2pt]
                                    e(\pi_R, g^{\alpha}) &= e(\pi_R', g), \\[-2pt]
                                    e(\pi_O, g^{\alpha}) &= e(\pi_O', g), \\[-2pt]
                                    e(\pi_H, g^{\alpha}) &= e(\pi_H', g).
                                \end{align*}
                                \vspace{-16pt}
                            }}
                            \onslide<8->{\item $e(\pi_L, g^{\gamma\beta_L}) \cdot e(\pi_R, g^{\gamma\beta_R}) \cdot \\ \cdot e(\pi_O, g^{\gamma\beta_O}) = e(\pi_{\beta}, g^{\gamma})$}
                        \end{itemize}}
                    \end{minipage}
                };}
            \end{tikzpicture}
        \end{center}
    \end{frame}

    \section{Real Protocols}

    \begin{frame}{Public and Private Inputs}
        So far, the verifier did not use the public statement at all. Split the indices of the extended witness $\mathbf{z}$ into $\mathcal{I}_{\text{io}}$ (input/output, public) and $\mathcal{I}_{\text{mid}}$ (intermediate, private), with $\mathcal{I}_{\text{mid}} \cap \mathcal{I}_{\text{io}} = \emptyset$, $\mathcal{I}_{\text{mid}} \cup \mathcal{I}_{\text{io}} = \{1,\dots,n\}$.
        \pause

        \begin{equation*}
            \pi_L = \underbrace{\textcolor{oc-pink-7}{g^{\sum_{i \in \mathcal{I}_{\text{mid}}} z_i\ell_i(\tau)}}}_{\textcolor{oc-pink-7}{\text{new $\pi_L$}}}
            \times
            \underbrace{\chg{g^{\sum_{i \in \mathcal{I}_{\text{io}}} z_i\ell_i(\tau)}}}_{\chg{\pi_{\ell,\text{io}}}}.
        \end{equation*}
        \pause

        \begin{remark}
            The prover computes $\pi_L,\pi_R,\pi_O$ only over $\mathcal{I}_{\text{mid}}$, while the verifier restores $\pi_L^* \gets \pi_L\prod_{i \in \mathcal{I}_{\text{io}}} (g^{\ell_i(\tau)})^{z_i}$ (and similarly $\pi_R^*$, $\pi_O^*$) from the public values and uses them in the polynomial equality check.
        \end{remark}
    \end{frame}

    \begin{frame}{Complexity of the Basic Protocol}
        \begin{proposition}[Complexity of the Basic Protocol]
            Suppose circuit consists of $n$ gates. Then, the complexity of the basic protocol is as follows:
            \begin{itemize}
                \item \textbf{Proof Size:} $O(1)$ --- constant number of group elements.
                \item \textbf{Setup Time:} $O(n)$ --- calculating powers of $\tau$, evaluations at $\tau$.
                \item \textbf{Prover Time:} $O(n \log n)$ --- using FFT and wise choice of $\Omega$.
                \item \textbf{Verifier Time:} $O(1)$ --- constant number of pairings.
            \end{itemize}
            \pause

            However, $O(1)$ is not very descriptive for proof and verifier complexities, so let us provide a more detailed analysis.
            \begin{itemize}
                \item \textbf{Proof Size:} 9 $\mathbb{G}$ group elements.
                \item \textbf{Verifier Time:} 15 pairings.
            \end{itemize}
        \end{proposition}
        \pause

        We can do better!
    \end{frame}

    \begin{frame}{Pinocchio Protocol}
        \begin{block}{Idea}
            In toxic waste, include $\rho_L,\rho_R \leftarrowS \mathbb{F}$, set $\rho_O \gets \rho_L\rho_R$, and define the following generators:
            \begin{equation*}
                g_L \gets g^{\rho_L}, \quad g_R \gets g^{\rho_R}, \quad g_O \gets g^{\rho_O}
            \end{equation*}
        \end{block}
        \pause

        \begin{remarkblock}{Reason}
            Such choice of generators reduces 15 pairings to \textbf{11 pairings}. Additionally, we have only \textbf{8 group elements} in the proof.
        \end{remarkblock}
    \end{frame}

    \begin{frame}{Groth16 Protocol}
        \begin{block}{Idea: Generic Group Model}
            Use \textbf{Generic Group Model} (GGM) technique. Simply put, GGM allows the adversary to only make \textbf{oracle requests} to compute the group operations. For example, having a set $\{g^{\alpha r_i(\tau)}\}_{i \in [d]}$, adversary can compute only linear combinations of these values. In the particular case of Groth16, instead of considering $\ell_i(T)$, $r_i(T)$, and $o_i(T)$ separately, we construct their linear combinations as $Q_i(T) := \beta \ell_i(T) + \alpha r_i(T) + o_i(T)$, where $\alpha$ and $\beta$ are toxic parameters.
        \end{block}
        \pause

        \begin{remarkblock}{Asymmetric Pairing}
            Groth16 uses the asymmetric pairing $e: \mathbb{G}_1 \times \mathbb{G}_2 \to \mathbb{G}_T$ with generators $g_1 \in \mathbb{G}_1$ and $g_2 \in \mathbb{G}_2$: it is more efficient than the symmetric one.
        \end{remarkblock}
    \end{frame}

    \begin{frame}{Groth16 Protocol: Setup Procedure}
        Sample $\alpha,\beta,\gamma,\delta,\tau \leftarrowS \mathbb{F}$. The proving key is formed as follows:
        \begin{align*}
            \mathsf{pp} \gets \Big(&g_1^{\alpha},g_1^{\beta},g_1^{\delta},\left\{g_1^{\tau^i},g_1^{(\beta \ell_i(\tau) + \alpha r_i(\tau) + o_i(\tau))/\gamma}, g_1^{\tau^i z(\tau)/\delta}\right\}_{i \in [n]}, \\
            &g_2^{\beta}, g_2^{\delta}, g_2^{\gamma}, \{g_2^{\tau^i}\}_{i \in [d]}\Big)
        \end{align*}
    \end{frame}

    \begin{frame}{Groth16 Protocol: Proving Procedure}
        Sample random $\delta_L,\delta_R \leftarrowS \mathbb{F}$ and compute the following values:
        \begin{gather*}
            \pi_L \gets g_1^{\alpha + \sum_{i=1}^n z_i\ell_i(\tau) + \delta_L\delta}, \quad \pi_R \gets g_2^{\beta + \sum_{i=1}^n z_ir_i(\tau) + \delta_R\delta}, \\ \pi_O \gets g_1^{\frac{Q_{\text{mid}}(\tau)+h(\tau)z(\tau)}{\delta} + \ell\delta_R + r\delta_L - \delta_L\delta_R\delta},
        \end{gather*}
        \pause

        where by $Q_{\text{mid}}$ we denoted the following expression:
        \begin{equation*}
            Q_{\text{mid}}(\tau) = \sum_{i \in \mathcal{I}_{\text{mid}}} z_iQ_i(\tau) = \sum_{i \in \mathcal{I}_{\text{mid}}}z_i(\beta \ell_i(\tau) + \alpha r_i(\tau) + o_i(\tau))
        \end{equation*}
    \end{frame}

    \begin{frame}{Groth16 Protocol: Verification Procedure}
        The verifier first calculates the following value:
        \begin{equation*}
            \pi_{\text{io}} \gets g_1^{\sum_{i \in \mathcal{I}_{\text{io}}}z_i(\beta \ell_i(\tau) + \alpha r_i(\tau) + o_i(\tau))/\gamma},
        \end{equation*}
        \pause

        and then checks the following single condition:
        \begin{equation*}
            e(\pi_L, \pi_R) = e(g_1^{\alpha}, g_2^{\beta})e(\pi_{\text{io}},g_2^{\gamma})e(\pi_O,g_2^{\delta})
        \end{equation*}
        \pause

        \begin{remarkblock}{Note}
            $e(g_1^{\alpha}, g_2^{\beta})$ can be additionally hard-coded in the verifier, thus reducing the number of pairings to 3. Finally, the proof's size is now reduced to 3 group elements: two from $\mathbb{G}_1$, and one from $\mathbb{G}_2$.
        \end{remarkblock}
    \end{frame}

    \makeclosingframe
\end{document}
